\section{The microscopic posterior of the arithmetic endpoint classes}
\label{sec:v42-phase-posterior}
The section residues are physical source masses. Their role can be
made more explicit than a scalar correction to the local theorem:
they determine a finite conditional law of the two endpoint phases.
The phase-class indicators need not be strong multipliers. We instead
realize the required initial densities as finite sums of physical
peripheral vectors. This keeps the exact conditioning event unchanged.

Fix $R\in I$ in this section. Use the generator $q_*$, order $d=d_R$,
phase masses $w_l$, and residue $r=r_R(k,n,m)$ of
\eqref{eq:v40-section-phase-masses}--\eqref{eq:v40-arithmetic-factor}.
Let $\omega=e^{2\pi i/d}$ and define $H_R(x)\in\Z/d\Z$ by
$q_*(x)=\omega^{H_R(x)}$ on its invariant full-measure set. Write
\[
 S_r=\sum_{l\bmod d}w_lw_{l+r},\qquad
 E=\{K_{n,R}=k,\ N_{n,R}=m,\ T_{n,R}-t\in J\},
\]
where $J$ is a fixed bounded interval of positive length. On $E$ the
exact phase equation gives
\begin{equation}\label{eq:v42-exact-phase-pair}
 H_R((F_R^*)^n x)=H_R(x)+r\pmod d.
\end{equation}
For $S_r>0$ put
\begin{equation}\label{eq:v42-phase-posterior}
 \pi_{R,r}(l,l')=
   \one_{\{l'=l+r\}}\frac{w_lw_{l+r}}{S_r}.
\end{equation}
This is a probability, invariant under simultaneous relabeling of the
phase classes. For $d=1$ it is the point mass at $(0,0)$.

\begin{lemma}[Strong initial densities for measurable phase classes]
\label{lem:v42-class-source}
Each nonnegative bounded density
$a_l=\eta_R\one_{\{H_R=l\}}$ has $a_l\nu\in\mathcal B$ on the
local collision space at this fixed radius. Its construction does not
assert that multiplication by $q_*$ or by the phase-class indicator is
bounded on $\mathcal B$.
\end{lemma}
\begin{proof}
The full peripheral theorem gives a strong physical eigenvector
$q_*^j\nu$ for every $0\le j<d$. Indeed $q_*^j$ solves the phase
equation for the $j$th resonance; its bounded physical representative
agrees up to a constant with the strong eigenvector by simplicity.
The exact section multiplier is bounded on the strong space. Therefore
\begin{equation}\label{eq:v42-class-source-fourier}
 a_l\nu=\frac1d\sum_{j=0}^{d-1}
                        \omega^{-jl}M_{\eta_R}(q_*^j\nu)
\end{equation}
is a strong vector. The identity is finite Fourier inversion of the
pointwise relation $q_*^d=1$. Its physical representative is the
nonnegative function $a_l$. Only multiplication by $\eta_R$, not by
an unknown phase, has been used.
\end{proof}

\begin{theorem}[Exact-index arithmetic endpoint posterior]
\label{thm:v42-endpoint-posterior}
As $m\to\infty$, uniformly over the central targets $|Z|\le M$ at
this fixed radius and over $l\in\Z/d\Z$,
\begin{equation}\label{eq:v42-class-local-law}
 m^2\nu_R^*(E\cap\{H_R=l\})
       =\frac d c |J|g_{\Omega_R}(Z)w_lw_{l+r}+o_R(1).
\end{equation}
If $w_lw_{l+r}=0$, the event on the left is exactly null, not merely
asymptotically negligible. On every positive residue class $S_r>0$,
\begin{equation}\label{eq:v42-class-TV}
 d_{\rm TV}\!\left(
 \operatorname{Law}_{\nu_R^*(\cdot\mid E)}
       (H_R(x),H_R((F_R^*)^nx)),\ \pi_{R,r}\right)\longrightarrow0.
\end{equation}
Here $d_{\rm TV}$ on finite probability spaces is one half of the
sum of the absolute mass differences. Neither assertion is made
uniform through radius-dependent changes in the phase group.
\end{theorem}
\begin{proof}
For a compactly band-limited roof test, repeat the fixed-count inverse
in Theorem~\ref{thm:v40-arithmetic-return-LLT}, now with the strong
initial vector $a_l\nu$ and terminal multiplier $\eta_R$. This is a
legitimate endpoint pairing by Lemma~\ref{lem:v42-class-source}.
The general physical projection formula gives, at resonance $j$,
\[
 \nu(a_l\overline{q_*^j})\nu(\eta_Rq_*^j)
       =w_l\omega^{-jl}\sum_hw_h\omega^{jh}.
\]
The constant inverse phase is $\omega^{-jr}$. Hence its complete
finite residue sum is
\begin{equation}\label{eq:v42-class-residue-sum}
 \sum_{j=0}^{d-1}\omega^{-jr}
        w_l\omega^{-jl}\sum_hw_h\omega^{jh}
                          =d w_lw_{l+r}.
\end{equation}
All local curvatures are the same $\Omega_R$. At fixed $R$ there are
finitely many initial vectors and resonances, so the Gaussian
localization errors are uniform in $l$ and in central targets. This
proves \eqref{eq:v42-class-local-law} for the band-limited test with
$|J|$ replaced by its integral. To pass to $J$, multiply ordered
Fourier envelopes by the positive physical source $a_l(x)\eta_R(T^mx)$
and the unchanged exact lattice/occupation indicators. Fix the roof
band, take $m\to\infty$, and then enlarge the band. All roof resonance
coordinates vanish, so both envelopes have the same coefficient
\eqref{eq:v42-class-residue-sum}. This proves the stated interval law.

Equation~\eqref{eq:v42-exact-phase-pair} confines the endpoints to the
pair $(l,l+r)$. If its initial or terminal section class is null,
invariance of the collision measure and its fixed normalization on the
section make the event null. For $S_r>0$, summing
\eqref{eq:v42-class-local-law} in $l$ gives the denominator
$(d/c)|J|g_{\Omega_R}(Z)S_r+o_R(1)$. It is bounded away from zero
on the scaled central sets. Division proves convergence of every
finite conditional mass. There are only finitely many positive $S_r$
at this fixed radius. Summing the finitely many errors proves
\eqref{eq:v42-class-TV} uniformly over them.
\end{proof}

\begin{corollary}[Phase classes and the actual Gaussian return bridge]
\label{cor:v42-phase-bridge-product}
Let $\mathcal Y_{n,R}$ be the return path tied to its actual endpoint
in \eqref{eq:v41-return-bridge}. On the same fixed-radius central
positive residue classes, its joint conditional law with the two phase
classes converges, in bounded-Lipschitz distance on the product of the
finite class space and continuous path space, to
\begin{equation}\label{eq:v42-phase-bridge-product}
        \pi_{R,r}\otimes\operatorname{Law}(\mathbb B_{D_R}).
\end{equation}
Thus the arithmetic posterior need not be uniform, but is
asymptotically independent of the pinned Gaussian fluctuations.
The statement uses a fixed positive roof interval, not conditioning
on a roof density value.
\end{corollary}
\begin{proof}
Insert the initial vector $a_l\nu$ in the chronological pinned
characteristic calculation of Lemma~\ref{lem:v41-pinned-factorization}.
At this fixed radius the resonances and their curvatures are fixed.
The block perturbations have zero weighted sum, so the drift and
Fourier/bridge cross terms cancel as before. The only changed
zeroth-order factor is \eqref{eq:v42-class-residue-sum}. Complementary
powers and projector errors obey the same estimates, with constants
allowed to depend on the finite list of $a_l\nu$. The sharp interval
error is bounded by the corresponding positive source-envelope
error before any probability division. Thus the conditional
finite-dimensional characteristic limit on each class with
$w_lw_{l+r}>0$ is the same Gaussian bridge.

The unnormalized fourth-moment bound of
Proposition~\ref{prop:v41-pinned-fourth} also holds after restricting
to $\{H_R=l\}$: its integrand is nonnegative and its source event
has only been restricted. The new class local law supplies its
positive denominator of order $m^{-2}$. This proves tightness on
each of the finitely many positive class pairs. The inverse occupation
clock and the exact compensation in
Theorem~\ref{thm:v41-actual-return-bridge} are pathwise identities
on this unchanged subevent, so the same return-clock conversion gives
$\mathbb B_{D_R}$. Class pairs of zero weight are exactly null.
Combining these conditional path limits with
\eqref{eq:v42-class-TV} proves the finite-mixture product limit.
\end{proof}

\paragraph{The arithmetic form of the raw target.}
The exact-index main term at fixed radius is the arithmetic expression
$c\,\mathfrak a_R(k,n,m)g_{\Omega_R}(Z)$, or equivalently
$\mathfrak a_R(k,n,m)g_{D_R}(V_n)$ in return normalization. The
radius-uniform object is the finite transition kernel
$\mathcal L_{m,R}$, not a pointwise choice of an exact resonance group.
These are the candidate main terms for a pointwise roof-density
extension of the corresponding interval theorems. The unmodulated
specialization requires the exact zero-residue criterion of
Theorem~\ref{thm:v41-zero-residue-criterion}.

Neither the finite posterior formula nor a finite packet average
proves that the concrete section phase masses equal $c/d$. The new
critical subtraction in Section~\ref{sec:v42-morse-packet} works with
each actual index and therefore does not average away this arithmetic.
Its boundary-source error must still be controlled pointwise before
the interval main term can be promoted to a raw density theorem.
